%%%%%%%%%%%%%%%%%%%%%%%%%%%%%%%%%%%%%%%%%%%%%%%%%%%%%%%%%%%%%
%% Begin exercise %%
%%%%%%%%%%%%%%%%%%%%%%%%%%%%%%%%%%%%%%%%%%%%%%%%%%%%%%%%%%%%%
\ex{Sensors}
% \newif
% \showsolutionstrue % <-- change to \showsolutionstrue to show

%%%%%%%%%%%%%%%%%%%%%%%%%%%%%%%%%%%%%%%%%%%%%%%%%%%%%%%%%%%%%
%% Task 2.1 %%
%%%%%%%%%%%%%%%%%%%%%%%%%%%%%%%%%%%%%%%%%%%%%%%%%%%%%%%%%%%%%

\task{Piezoelectric accelerometer under shock + drift}

You’re designing a 10 g full-scale accelerometer for impact testing using a PZT stack in the longitudinal mode (use d$_{33}$ = 593 pC/N). 
The proof mass is m = 5 g. The shock is a half-sine a(t) = a$_{\text{pk}}$ sin($\pi t / T$) with a$_{\text{pk}}$ = 500 m/s$^2$, T = 4 ms. 
The ceramic disk has active area A = 100 mm$^2$, thickness t = 0.5 mm, relative permittivity $\varepsilon_r$ = 1500. 
Assume quasi-static force transfer F(t) = m a(t) and a charge amplifier with C$_f$ = 100 pF, ideal op-amp.


\subtask{ Compute the sensor charge \( Q(t) \) and amplifier output \( v_o(t) \).}


\begin{solutionblock}
\noindent
Force:
\begin{align*}
    F(t) &= m\,a(t) = m\,a_{\text{pk}}\sin\!\left(\frac{\pi t}{T}\right).
\end{align*}

\noindent
Piezoelectric charge (direct effect, quasi-static transfer):
\begin{align*}
    Q(t) &= d_{33}\,F(t) = d_{33}\,m\,a_{\text{pk}}\sin\!\left(\frac{\pi t}{T}\right).
\end{align*}

\noindent
Charge amplifier ideal output:
\begin{align*}
    v_o(t) &= -\frac{Q(t)}{C_f}
            = -\frac{d_{33}\,m\,a_{\text{pk}}}{C_f}\sin\!\left(\frac{\pi t}{T}\right).
\end{align*}

\noindent
Calculating using the given values:-

\begin{align*}
    F_{\text{peak}} &= m\,a_{\text{pk}} = 0.005 \times 500 = 2.5~\mathrm{N},\\[6pt]
    Q_{\text{peak}} &= d_{33}\,F_{\text{peak}}
                    = 593\times10^{-12} \times 2.5
                    = 1.4825\times10^{-9}~\mathrm{C}
                    = 1482.5~\mathrm{pC},\\[6pt]
    v_{o,\text{peak}} &= -\frac{Q_{\text{peak}}}{C_f}
                      = -\frac{1.4825\times10^{-9}}{100\times10^{-12}}
                      = -14.825~\mathrm{V}.
\end{align*}

\noindent
Therefore, the time-domain results are:
\begin{align*}
    Q(t) &= 1.4825\times10^{-9}\sin\!\left(\frac{\pi t}{4~\text{ms}}\right)\ \text{C},\\[4pt]
    v_o(t) &= -14.825\sin\!\left(\frac{\pi t}{4~\text{ms}}\right)\ \text{V},
\end{align*}
for \( 0 \le t \le 4~\text{ms} \) (and zero outside the half-sine).

\bigskip
\noindent
\textit{Note:} The output \( v_o(t) \) follows a negative half-sine waveform with a peak amplitude of approximately \(-14.825~\text{V}\), due to the inverting nature of the charge amplifier.

\end{solutionblock}



\subtask{ Estimate the sensor capacitance \( C_s \) and discuss how \( C_s \) versus \( C_f \) affects the scale factor and low-frequency droop.}


\begin{solutionblock}
\textbf{Ideal charge amplifier:} For an ideal op-amp, the output depends only on the generated charge and the feedback capacitor:
\[
v_o(t) = -\frac{Q(t)}{C_f}
\]
so ideally the sensor capacitance \(C_s\) does not affect the scale factor, which is set by \(1/C_f\).
\[
C_s = \frac{\varepsilon_0 \varepsilon_r A}{t} \approx 2.7~\mathrm{nF}, \quad
\text{with } C_s \gg C_f, \]
\[\text{ the op-amp forces the signal charge onto } C_f:
\quad |v_o| = \frac{Q}{C_f}.\]
\text{Smaller } $C_f$ \text{ increases sensitivity but also raises noise and saturation risk.}

\end{solutionblock}



\subtask{ If the device warms by \( +40^{\circ}\mathrm{C} \), qualitatively discuss the expected sensitivity shift and explain why a constant force cannot be measured indefinitely with a piezoelectric sensor.}


\begin{solutionblock}
Piezo sensors’ sensitivity depends on $d_{33}$ and $\varepsilon_r$, which decrease with temperature, 
while leakage and dielectric loss increase, causing faster signal droop. 
They generate charge under stress but cannot hold a DC output because leakage and dielectric 
absorption make the signal decay. 
Therefore, piezo sensors are best for measuring dynamic or transient forces, not constant ones.
\end{solutionblock}


%%%%%%%%%%%%%%%%%%%%%%%%%%%%%%%%%%%%%%%%%%%%%%%%%%%%%%%%%%%%%
%% Task 2.2 %%
%%%%%%%%%%%%%%%%%%%%%%%%%%%%%%%%%%%%%%%%%%%%%%%%%%%%%%%%%%%%%

\task{Strain-gauge bridge}
A steel cantilever ($E = 210~\text{GPa}$) carries a tip load that produces surface strain $\varepsilon = 800~\mu\varepsilon$ at the gauge location.
Four identical $120~\Omega$ foil strain gauges (gauge factor $k \approx 2$) are arranged in a full bridge with $5.0~\text{V}$ excitation.
All four gauges experience the same temperature rise of $+30~^\circ\text{C}$.

\subtask{Derive the bridge output $U_A$ with two active gauges in tension and two in compression.}


\begin{solutionblock}
Take the standard Wheatstone full-bridge. Label resistances clockwise: 
$R_1$ (top left), $R_2$ (bottom left), $R_3$ (top right), $R_4$ (bottom right).

Choosing the following arrangement.
\[
R_1 = R + \Delta R \ (\text{tension}), \quad\\
R_2 = R - \Delta R \ (\text{compression}), \quad\\
R_3 = R - \Delta R \ (\text{compression}), \quad\\
R_4 = R + \Delta R \ (\text{tension})
\]

Node voltages (top node = $U_\text{exc}$, bottom node = 0):
\[
V_A = U_\text{exc} \frac{R_2}{R_1 + R_2}, \qquad
V_B = U_\text{exc} \frac{R_4}{R_3 + R_4}
\]

Bridge output:
\[
U_A = V_A - V_B
\]

Substitute the resistances and linearize for $\Delta R \ll R$:
\[
U_A \approx -\,U_\text{exc} \frac{\Delta R}{R} = -\,U_\text{exc}\, k\, \varepsilon
\]

(Sign depends on which midpoint is subtracted from which; magnitude is $U_\text{exc} k \varepsilon$.)

Numerical value:
\[
U_A = U_\text{exc}\, k\, \varepsilon = 5.0 \times 2 \times 800 \times 10^{-6} = 0.008~\text{V} = 8.0~\text{mV}
\]

Thus, the full-bridge gives an 8 mV output for the stated strain.

\end{solutionblock}

\subtask{ Show why a full bridge largely cancels temperature effects.}  
\begin{solutionblock}
\textbf{Temperature Compensation in a Full Bridge:}  

If all four gauges experience the same temperature change, adding $\Delta R_T$ to each arm, the common-mode shifts cancel out. The bridge output $V_A - V_B$ depends only on strain.  

Practical note: small errors can occur due to TCR mismatches, uneven mounting, or self-heating, but a full bridge greatly reduces temperature effects compared to a single gauge or half-bridge.
\end{solutionblock}


\subtask{If the adhesive allows only $I_\text{exc} \le 15~\text{mA}$ to avoid self-heating, check compliance.}


\begin{solutionblock}
For a full bridge with \(U_\text{exc} = 5\ \text{V}\) and two series legs (\(\approx 2R\)) in parallel:

\[
R_\text{eq} = \frac{(R_1+R_2)(R_3+R_4)}{(R_1+R_2)+(R_3+R_4)} \approx R
\]

\[
I_\text{exc} = \frac{U_\text{exc}}{R_\text{eq}} \approx \frac{5}{120} = 41.7\ \text{mA} > 15\ \text{mA} \ (\text{too high})
\]

Maximum excitation for \(I_\text{exc} \le 15\ \text{mA}\):

\[
U_{\text{exc,max}} = I_\text{max} R = 0.015 \times 120 = 1.8\ \text{V}
\]

Power at \(U_\text{exc} = 5\ \text{V}\):

\[
P_\text{tot} = \frac{U_\text{exc}^2}{R_\text{eq}} = 0.208\ \text{W}, \quad P_\text{gauge} \approx 52\ \text{mW}
\]

At \(U_\text{exc} = 1.8\ \text{V}\), \(P_\text{gauge} \approx 6.75\ \text{mW}\) (safe).

\end{solutionblock}

%%%%%%%%%%%%%%%%%%%%%%%%%%%%%%%%%%%%%%%%%%%%%%%%%%%%%%%%%%%%%
%% Task 2.3 %%
%%%%%%%%%%%%%%%%%%%%%%%%%%%%%%%%%%%%%%%%%%%%%%%%%%%%%%%%%%%%%

\task{Silicon piezoresistive membrane pressure sensor: bias-point design}
A square silicon diaphragm carries four diffused piezoresistors in a full bridge. At rated pressure, the longitudinal surface strain is \(+500~\mu\varepsilon\) at two opposite edges and \(-500~\mu\varepsilon\) at the other two. Each resistor is nominally \(3~\text{k}\Omega\), with piezoresistive gauge factor \(k_\pi = 80\) (n-type).

\subtask{Find the small-signal bridge sensitivity \(S = dU_A/dp\)}


\begin{solutionblock}
For each resistor, \(\Delta R / R = k_\pi \varepsilon\). With two \(+\varepsilon\) and two \(-\varepsilon\):

\[
\frac{U_A}{U_\text{exc}} \approx \frac{k_\pi \varepsilon}{2}
\]

With \(k_\pi = 80\) and \(\varepsilon = 500\times10^{-6}\):

\[
U_A \approx 0.02\, U_\text{exc} \ (\text{20 mV/V})
\]

Bridge sensitivity formula:

\[
S = \frac{dU_A}{dp} = U_\text{exc} \, k_\pi \, \frac{d\varepsilon}{dp}
\]
\end{solutionblock}


\subtask{Discuss why silicon (vs. metal SG) achieves much higher sensitivity and how temperature-dependent \(\rho\) enters}


\begin{solutionblock}
Silicon’s large piezoresistive effect (\(k_\pi \approx 80\)) gives much higher sensitivity than metal gauges (\(k \approx 2\)). 
Temperature affects resistivity \(\rho(T)\), adding a TCR term to \(\Delta R/R\).
\[
\frac{\Delta R}{R} = k_\pi \, \varepsilon + \alpha_T \, \Delta T
\]
The bridge cancels common temperature shifts, but additional compensation is usually required.
\end{solutionblock}


%%%%%%%%%%%%%%%%%%%%%%%%%%%%%%%%%%%%%%%%%%%%%%%%%%%%%%%%%%%%%
%% Task 2.4 %%
%%%%%%%%%%%%%%%%%%%%%%%%%%%%%%%%%%%%%%%%%%%%%%%%%%%%%%%%%%%%%

\task{Hall vs field-plate (Gaussian/MR) for isolated current sensing}
Measure up to 200A DC through a busbar with 1~mm air gap using either:\\
A. a Hall plate (\(d = 10~\mu\text{m},~ b = 200~\mu\text{m},~ I_x = 5~\text{mA}\)); or\\ 
B. a field-plate MR element (\(R_0 = 1~\text{k}\Omega,~ \alpha \approx -0.004~\text{K}^{-1}\)).\\  
Flux density at full scale: \(B = 80~\text{mT}\).

\subtask{ For A, estimate \(U_H\) using \(U_H = (A_H B/d) I_x\) with a typical semiconductor \(A_H\).}


\begin{solutionblock}
\[
U_H = \frac{A_H B}{d} I_x
\]
With \(A_H = 5\times10^{-4}\ \text{m}^3/\text{C}\), \(B = 0.08~\text{T}\), \(d = 10^{-5}~\text{m}\), \(I_x = 5~\text{mA}\):
\[
U_H = \frac{5\times10^{-4}\cdot0.08}{1\times10^{-5}}\cdot0.005 = 0.02~\text{V} = 20~\text{mV}
\]
Typical Hall output = 4–40~mV for \(A_H\) between \(10^{-4}\) and \(10^{-3}\ \text{m}^3/\text{C}\)
\end{solutionblock}


\subtask{ For B, use 5~mA constant-current bias and the \(R(B)\) characteristic to estimate \(\Delta V\).}


\begin{solutionblock}
With MR ratio \(\Delta R / R_0 \approx 0.02\) at \(B = 80~\text{mT}\):
\[
\Delta V = I\,R_0\,\frac{\Delta R}{R_0} = 0.005 \times 1000 \times 0.02 = 0.1~\text{V} = 100~\text{mV}
\]
So MR yields = 100~mV output, higher than the Hall signal.
\end{solutionblock}


\subtask{ Compare temperature robustness and reasons to pick A vs. B.}


\begin{solutionblock}
\textbf{Temperature sensitivity:}  
MR has \(\alpha \approx -0.004~\text{K}^{-1}\), giving:
\[
\frac{\Delta R}{R}\bigg|_T = \alpha \Delta T,\quad \Delta V_T = I\,R_0\,\alpha \Delta T \approx -20~\text{mV/K}
\]
$\Rightarrow$ strong temperature drift, requires compensation.

Hall devices also vary with temperature, but typically only a few \% over tens of °C due to changes in \(A_H\) and mobility — smaller drift and often compensated in integrated Hall ICs.

\textbf{Choice:}  
Hall $\rightarrow$ simple, linear DC response, moderate sensitivity (~20~mV), better thermal stability.  
MR $\rightarrow$ higher sensitivity (~100~mV) but large \(|\alpha|\), needs calibration or temperature control.
\end{solutionblock}


%%%%%%%%%%%%%%%%%%%%%%%%%%%%%%%%%%%%%%%%%%%%%%%%%%%%%%%%%%%%%
%% Task 2.5 %%
%%%%%%%%%%%%%%%%%%%%%%%%%%%%%%%%%%%%%%%%%%%%%%%%%%%%%%%%%%%%%

\task{Magnetostrictive position sensor: time-of-flight and linearity}

A magnetostrictive rod (wave speed $v = 2850~\text{m/s}$) is used in a non-contact absolute position transducer. A torsional pulse is launched at $t = 0$ and detected after interacting with a magnet placed at unknown position $x$ along a $0.6~\text{m}$ rod. The measured time is $t = 152~\mu\text{s}$.

\subtask{ Estimate $x$.}


\begin{solutionblock}
Assuming the pulse travels one way (no reflection path),
\[
x = v t = 2850 \times 152 \times 10^{-6} 
   = 0.4332~\text{m} \approx 433~\text{mm.}
\]
\end{solutionblock}


\subtask{ Correct for the electronics delay of $3.0~\mu\text{s}$.}


\begin{solutionblock}
Corrected travel time:
\[
t_{\text{true}} = t - 3.0~\mu\text{s} = 152 - 3 = 149~\mu\text{s.}
\]
Then,
\[
x_{\text{corr}} = v t_{\text{true}} 
= 2850 \times 149 \times 10^{-6} 
= 0.4247~\text{m} \approx 425~\text{mm.}
\]
The delay caused an overestimate of $\approx 8.5~\text{mm.}$
\end{solutionblock}


\subtask{ Temperature effect on position error.}


\begin{solutionblock}
Given $ \frac{\Delta v}{v} = -0.02\%\!/\text{K}$, for a $+20~\text{K}$ change:
\[
\frac{\Delta v}{v} = -0.0002 \times 20 = -0.004.
\]
Since $x \propto v$, 
\[
\Delta x = -0.004 \times 0.4247 = -1.70\times10^{-3}~\text{m} 
\approx -1.7~\text{mm.}
\]
Thus, a $+20~\text{K}$ temperature rise (uncorrected) causes $\approx -1.7~\text{mm}$ error.
\end{solutionblock}


\subtask{ Principle of operation (Villari and Wiedemann effects).}


\begin{solutionblock}
\textbf{Wiedemann effect (launch):} A current pulse through the rod interacts with a magnetic bias, producing a torsional strain wave. \\
\textbf{Villari effect (detection):} The torsional stress under the external magnet alters the local magnetization, inducing a voltage in a pickup coil. \\
The time-of-flight of this torsional wave gives the absolute position of the magnet.
\end{solutionblock}




%%%%%%%%%%%%%%%%%%%%%%%%%%%%%%%%%%%%%%%%%%%%%%%%%%%%%%%%%%%%%
%% Task 2.6 %%
%%%%%%%%%%%%%%%%%%%%%%%%%%%%%%%%%%%%%%%%%%%%%%%%%%%%%%%%%%%%%

\task{Eddy-current proximity sensor: oscillator detuning and metal target}

An inductive proximity sensor uses an LC oscillator with $L = 120~\mu\text{H}$, $C = 680~\text{pF}$.  
A conductive aluminium target causes an inductance drop $\Delta L = -12\%$ at the switch point due to eddy currents.

\subtask{ Find the free-run frequency and the switch-point frequency.}


\begin{solutionblock}
The free-run oscillator frequency is
\[
f_0 = \frac{1}{2\pi\sqrt{LC}}
     = \frac{1}{2\pi\sqrt{120\times10^{-6}\times680\times10^{-12}}}
     \approx 557~\text{kHz.}
\]

At the switch point, $L' = 0.88L$:
\[
f_s = \frac{1}{2\pi\sqrt{L'C}}
     = \frac{f_0}{\sqrt{0.88}}
     \approx 594~\text{kHz.}
\]
\end{solutionblock}


\subtask{ Explain why conductive non-magnetic targets reduce $L$ and how coil $Q$ plays in detection margin.}


\begin{solutionblock}
Conductive targets (e.g., aluminium) support eddy currents that \emph{oppose} the coil’s magnetic flux (Lenz’s law), thus reducing effective $L$ and increasing oscillator frequency.

Coil quality factor:
\[
Q = \frac{\omega L}{R} = \frac{2\pi f_0 L}{R}.
\]
High $Q$ means a sharp, strong oscillation and higher sensitivity to $\Delta L$; low $Q$ broadens resonance and reduces detection margin.

Example: for $L=120~\mu\text{H}$, $R=10~\Omega$, $f_0=557~\text{kHz}$,
\[
Q = \frac{2\pi(5.57\times10^5)(1.2\times10^{-4})}{10} \approx 4.2,
\]
a typical moderate value for practical proximity sensors.
 
Eddy currents $\Rightarrow$ lower $L$, higher $f$.\\
High $Q$ $\Rightarrow$ sharper frequency or amplitude change, improving switch detection.
\end{solutionblock}



%%%%%%%%%%%%%%%%%%%%%%%%%%%%%%%%%%%%%%%%%%%%%%%%%%%%%%%%%%%%%
%% Task 2.7 %%
%%%%%%%%%%%%%%%%%%%%%%%%%%%%%%%%%%%%%%%%%%%%%%%%%%%%%%%%%%%%%

\task{Capacitive thickness gauge with fringing}

A non-contact thickness sensor measures a plastic sheet ($\varepsilon_r = 3.0$) with parallel electrodes $A = 4~\text{cm}^2$ and nominal gap $d = 0.50~\text{mm}$.  
Sheet thickness variations $\Delta d$ are to be resolved.

\subtask{ Compute nominal capacitance $C_0$.}


\begin{solutionblock}
\[
C_0 = \varepsilon_0 \varepsilon_r \frac{A}{d}
     = \frac{8.854\times10^{-12} \cdot 3.0 \cdot 4.0\times10^{-4}}{5.0\times10^{-4}}
     = 2.12\times10^{-11}~\text{F} = 21.2~\text{pF.}
\]
\end{solutionblock}


\subtask{ If the sheet bows, causing a +5\% fringing-equivalent increase in effective area, estimate the apparent thickness error if the parallel-plate model is used.}


\begin{solutionblock}
If $A_{\rm eff} = 1.05A$, then $C_{\rm meas} = 1.05C_0$ while the model assumes $A$ constant:
\[
d_{\rm inf} = \frac{\varepsilon_0 \varepsilon_r A}{C_{\rm meas}}
             = \frac{d}{1.05} \Rightarrow
\frac{\Delta d}{d} = \frac{1}{1.05} - 1 \approx -4.76\%.
\]
Numerically, $\Delta d \approx -0.024~\text{mm} = -23.8~\mu\text{m}$.  
Thus, ignoring fringing makes the sheet appear $\approx$24~µm thinner.
\end{solutionblock}


%%%%%%%%%%%%%%%%%%%%%%%%%%%%%%%%%%%%%%%%%%%%%%%%%%%%%%%%%%%%%
%% Task 2.8 %%
%%%%%%%%%%%%%%%%%%%%%%%%%%%%%%%%%%%%%%%%%%%%%%%%%%%%%%%%%%%%%

\task{Field-plate (Gaussian effect) tacho vs toothed wheel: waveform and calibration}

A field-plate bridge measures speed via a steel toothed wheel. Bridge excitation \(U_{\rm exc}=3.3\ \text{V}\), nominal arm \(R_0=1.2\ \text{k}\Omega\). At the peak of a tooth the active element resistance increases by \(+8\%\).

\subtask{ Output waveform and shape}


\begin{solutionblock}
As each tooth passes, the magnetic field \(B\) rises and falls, changing the resistance as \(R = R_0(1 \pm \delta)\).  
The bridge output is:
\[
U_A = U_{\rm exc}\,\delta.
\]
Because \(R(B)\) is nonlinear and the field shape around the tooth is not perfectly sinusoidal, the output is a \emph{quasi-sinusoidal but slightly distorted} waveform.
\end{solutionblock}


\subtask{ Estimate peak-to-peak output}


\begin{solutionblock}
With opposite arms varying by \(\pm8\%\):
\[
\Delta U_A \approx 2\,U_{\rm exc}\,\delta_{\max} = 2 \times 3.3 \times 0.08 \approx 0.53~\text{V}_{\rm p-p}.
\]
Hence about \(0.26~\text{V}\) amplitude on each side of zero.
\end{solutionblock}


\subtask{ If the element self-heats by $+15~\mathrm{K}$, estimate the DC drift using $\alpha \approx -0.004~\mathrm{K}^{-1}$ and discuss compensation.
}


\begin{solutionblock}
Self-heating by \(+15~\text{K}\) and \(\alpha = -0.004~\text{K}^{-1}\) gives:
\[
\frac{\Delta R}{R} = \alpha \Delta T = -0.004 \times 15 = -0.06,
\]
i.e. a \(-6\%\) baseline drop, causing a DC offset in the bridge.  
Use ratiometric or differential sensing and temperature correction to cancel this drift.
\end{solutionblock}


%%%%%%%%%%%%%%%%%%%%%%%%%%%%%%%%%%%%%%%%%%%%%%%%%%%%%%%%%%%%%
%% Task 2.9 %%
%%%%%%%%%%%%%%%%%%%%%%%%%%%%%%%%%%%%%%%%%%%%%%%%%%%%%%%%%%%%%

\task{Thermistor/RTD Self-Heating: Measurement-Current Limit}

An RTD with $R(25^\circ \text{C}) = 100~\Omega$ measures air temperature. The thermal resistance to ambient is $\Theta = 400~\text{K/W}$. The readout uses a constant current $I$. Allow at most $+0.2~^\circ\text{C}$ self-heating error.

\subtask{Find the maximum measurement current $I_{\rm max}$.}


\begin{solutionblock}
The self-heating temperature rise is

\[
\Delta T = \Theta I^2 R.
\]

Solving for $I_{\rm max}$:

\[
I_{\rm max} = \sqrt{\frac{\Delta T}{\Theta R}}
             = \sqrt{\frac{0.2}{400 \cdot 100}}
             \approx 2.24~\text{mA}.
\]
\end{solutionblock}


\subtask{If the response time is dominated by thermal capacitance $C_{\rm th} = 30~\text{mJ/K}$, estimate the time constant $\tau$ and discuss the speed/accuracy trade-off.}


\begin{solutionblock}
The thermal time constant is

\[
\tau = \Theta \cdot C_{\rm th} \approx 400 \cdot 0.03 = 12~\text{s}.
\]

\textit{Note:} Increasing $I$ improves the electrical signal-to-noise ratio (SNR), but it also increases self-heating $\Delta T$. A compromise must be made: choose $I$ such that the self-heating error is within the allowed budget while still meeting the required measurement bandwidth.
\end{solutionblock}



%%%%%%%%%%%%%%%%%%%%%%%%%%%%%%%%%%%%%%%%%%%%%%%%%%%%%%%%%%%%%
%% Task 2.10 %%
%%%%%%%%%%%%%%%%%%%%%%%%%%%%%%%%%%%%%%%%%%%%%%%%%%%%%%%%%%%%%

\task{Hall Current Sensor with Magnetic Circuit: Sizing and SNR}

A C-core with a $1.0~\text{mm}$ gap surrounds a conductor carrying up to $150~\text{A}$ DC. The mean magnetic path length is $l_m = 120~\text{mm}$, $\mu_r = 2000$. A Hall plate with thickness $d = 20~\mu\text{m}$ and bias current $I_x = 8~\text{mA}$ sits in the gap.

\subtask{Estimate the magnetic flux density $B$ in the gap at full scale.}


\begin{solutionblock}
The gap dominates the reluctance, so the magnetic flux density can be estimated as:

\[
B \approx \frac{\mu_0 I}{g}
\]

with $N = 1$ turn. Substituting values:

\[
B \approx \frac{4\pi \times 10^{-7} \cdot 150}{1.0\times 10^{-3}} \approx 0.19~\text{T}.
\]
\end{solutionblock}


\subtask{Estimate the Hall voltage $U_H$ for a semiconductor and a metal.}


\begin{solutionblock}
The Hall voltage is:

\[
U_H = \frac{A_H B}{d} I_x
\]

\begin{itemize}
    \item \textbf{Semiconductor:} $A_H$ is large $\rightarrow$ $U_H$ in practical range of tens of mV.
    \item \textbf{Metal:} $A_H$ is small $\rightarrow$ $U_H$ is tiny, often in $\mu$V range.
\end{itemize}
\end{solutionblock}

\subtask{Give two reasons why a Hall sensor is preferred over an inductive pickup in this application.}
\begin{solutionblock}
\begin{enumerate}
    \item Hall sensors measure \textbf{DC and low-frequency currents directly}, while inductive pickups require changing magnetic flux ($dB/dt$) and cannot measure steady DC.
    \item Hall sensors can provide a \textbf{compact, direct voltage output} proportional to current, simplifying signal processing, whereas inductive sensors need integration circuits for DC estimation.
\end{enumerate}
\end{solutionblock}


%%%%%%%%%%%%%%%%%%%%%%%%%%%%%%%%%%%%%%%%%%%%%%%%%%%%%%%%%%%%%
%% Task 2.11 %%
%%%%%%%%%%%%%%%%%%%%%%%%%%%%%%%%%%%%%%%%%%%%%%%%%%%%%%%%%%%%%

\task{Light Barrier Sensor for Fast Conveyor}

A sensor is needed for a "light barrier" application . The sensor must detect when a box on a fast-moving conveyor belt ($v = 2~\text{m/s}$) breaks a $1~\text{cm}$ wide light beam. The beam break time is:

\[
t = \frac{0.01~\text{m}}{2~\text{m/s}} = 5~\text{ms}.
\]

The sensor's output must react within this time. Two components are considered (both based on the internal photoelectric effect):

\begin{enumerate}
    \item Photoconductor (Photoresistor)  
    \item Silicon Photodiode
\end{enumerate}

The text states that photoconductors are slow (response in seconds at low light) while photodiodes are fast (usable up to GHz).  

\subtask{Explain the physical reason for this major difference in response speed.}


\begin{solutionblock}
\textbf{Photoconductor}: A photon excites an electron from the valence band to the conduction band, creating an electron-hole pair. 
These carriers increase the material's conductivity. 
The signal ends only when the carriers recombine, which is a slow, statistical process. 
At low light, the dark state is re-established very slowly (seconds).  

\textbf{Photodiode}: Electron-hole pairs are generated within or near the depletion zone of a pn-junction. 
The built-in electric field quickly sweeps electrons to the n-side and holes to the p-side, creating a photocurrent. 
The signal ends immediately when photons stop, as there are no free carriers waiting to recombine. 
This drift-current mechanism is vastly faster than the slow recombination in photoconductors.
\end{solutionblock}


\subtask{Which device is the only viable choice for this 5~ms application?}


\begin{solutionblock}
The application requires a response time $< 5~\text{ms}$:

\begin{itemize}
    \item The \textbf{Photoconductor} is not viable; its light-to-dark response is on the order of seconds.
    \item The \textbf{Photodiode} is suitable; its response time is limited by capacitance and is typically in the nanosecond-to-microsecond range, more than 1000 times faster than required.
\end{itemize}

Thus, the \textbf{Silicon Photodiode} is the correct choice.
\end{solutionblock}


%%%%%%%%%%%%%%%%%%%%%%%%%%%%%%%%%%%%%%%%%%%%%%%%%%%%%%%%%%%%%
%% Task 2.12 %%
%%%%%%%%%%%%%%%%%%%%%%%%%%%%%%%%%%%%%%%%%%%%%%%%%%%%%%%%%%%%%

\task{NTC Thermistor and Linearization}

An NTC thermistor (Negative Temperature Coefficient resistor) is used to measure temperature. Its resistance follows:

\[
R(T) = R_N \, e^{B \left( \frac{1}{T} - \frac{1}{T_N} \right)}
\]

with the following parameters:

\begin{itemize}
    \item Nominal resistance $R_N = 10~\text{k}\Omega$ at $T_N = 25^\circ\text{C}$ (298.15 K)
    \item Material constant $B = 3950~\text{K}$
\end{itemize}

\subtask{Calculate the sensor resistance $R(T)$ at $T = 0^\circ\text{C}$ and $T = 50^\circ\text{C}$.}


\begin{solutionblock}
\textbf{At $T = 0^\circ\text{C}$ (273.15 K):}

\[
\begin{aligned}
R(0^\circ\text{C}) &= 10~\text{k}\Omega \cdot e^{3950 \left(\frac{1}{273.15} - \frac{1}{298.15}\right)} \\
&= 10~\text{k}\Omega \cdot e^{3950 (0.003661 - 0.003354)} \\
&= 10~\text{k}\Omega \cdot e^{1.2126} \\
&\approx 33.6~\text{k}\Omega
\end{aligned}
\]

\textbf{At $T = 50^\circ\text{C}$ (323.15 K):}

\[
\begin{aligned}
R(50^\circ\text{C}) &= 10~\text{k}\Omega \cdot e^{3950 \left(\frac{1}{323.15} - \frac{1}{298.15}\right)} \\
&= 10~\text{k}\Omega \cdot e^{3950 (0.0030945 - 0.003354)} \\
&= 10~\text{k}\Omega \cdot e^{-1.025} \\
&\approx 3.59~\text{k}\Omega
\end{aligned}
\]

\textit{Self-check:} $R(0^\circ\text{C}) > R_N > R(50^\circ\text{C})$, consistent with NTC behavior.
\end{solutionblock}


\subtask{Explain qualitatively how adding a parallel resistor $R_P$ helps linearize the sensor output.}


\begin{solutionblock}
The $R_T(T)$ curve of the NTC is convex (steep at low temperatures, flatter at high temperatures). Adding a parallel resistor $R_P$ gives an equivalent resistance:

\[
R_{\rm eq} = \frac{R_T \cdot R_P}{R_T + R_P}
\]

\begin{itemize}
    \item \textbf{At low temperatures (0°C):} $R_T = 33.6~\text{k}\Omega \gg R_P$. Then $R_{\rm eq} \approx R_P = 10~\text{k}\Omega$, reducing the effective slope.
    \item \textbf{At high temperatures (50°C):} $R_T = 3.59~\text{k}\Omega \ll R_P$. Then $R_{\rm eq} \approx R_T = 3.59~\text{k}\Omega$, leaving the slope almost unchanged.
\end{itemize}

The parallel resistor compresses the high-resistance (low-temperature) end of the curve more than the low-resistance (high-temperature) end, flattening the exponential behavior and making the response more linear over the 0–50°C range. This reduces sensitivity slightly but significantly improves linearity for simple analog voltage dividers.
\end{solutionblock}


\end{document}
